\documentclass{article}

% Packages
\usepackage{geometry}
\usepackage{graphicx}
\usepackage{lipsum}
\usepackage[portuguese]{babel}

% Page setup
\geometry{a4paper, margin=2cm}

\begin{document}

% cover
\thispagestyle{empty} % Remove page number from first page

\begin{figure}[t]
    \includegraphics[width=3cm]{images/logo-puc-minas.png}
    \hspace{0.02\textwidth}
    \vline%
    \hspace{0.04\textwidth}
    \includegraphics[width=3cm]{images/logo-icei.jpeg}
\end{figure}

\hrulefill%
\vspace{\baselineskip}

\Large\noindent
\textbf{Pontifícia Universidade Católica de Minas Gerais} \\
\textbf{Instituto de Ciências Exatas e Informática} \\
\textbf{Departamento de Engenharia de Computação}

\begin{center}
    \vfill
    \Huge\textbf{Relatório: Trabalho Prático 2} \\
    \vspace{0.5cm}
    \Large\textbf{Registradores em VHDL} \\
    \vspace{1cm}
    \large \textbf{Professor(es)}: Antônio Hamilton Magalhães\\
    \vspace{0.5cm}
    \large \textbf{Aluno(s)}: Bruno Luiz Dias Alves de Castro \\
    \large \hspace{0.75cm} Rafael Ramos de Andrade \\
    \vfill
    \large Belo Horizonte \\ Campus Coração Eucarístico \\
    \vspace{\baselineskip}
    \large \today
\end{center}

% table of contents
\newpage
\thispagestyle{empty}
\tableofcontents

% body
\newpage
\large % document text size

\section{Introdução}

Durante as aulas da disciplina de Sistemas Reconfiguráveis, fomos introduzidos à linguagem VHDL. VHDL (\textbf{V}HSIC \textbf{H}ardware \textbf{D}escription \textbf{L}anguage) é uma linguagem de descrição de hardware. Com ela, podemos montar circuitos lógicos de maneira totalmente textual, o que garante à linguagem uma grande vantagem ante à soluções visuais.

\subsection{TP1}

Como primeiro trabalho prático (TP1), são propostas as montagens de dois circuitos utilizando a linguagem VHDL para a placa de desenvolvimento Altera:

\begin{itemize}
    \item Um ADDR\_MUX (Multiplexador de Endereçamento) (Secção 2)
    \item Uma ULA (Unidade Lógica Aritimética) (Secção 3)
\end{itemize}

Ambos os circuitos devem ser desenvolvidos utilizando programação concorrente, ou seja, sem trechos sequênciais no código-fonte.

\subsubsection{Objetivos}

Entre os objetivos que temos com o desenvolvimento deste trabalho prático podemos listar:

\begin{itemize}
    \item Aprender conceitos básicos da linguagem VHDL;
    \item Implementar utilizando programação concorrente os dois circuitos propostos;
    \item Compilar os circuitos e testar os resultados na placa de desenvolvimento Altera;
\end{itemize}

\section{ADDR\_MUX}
\lipsum[1]

\section{ULA (Unidade Lógica Aritimética)}

A ULA (\textbf{U}nidade \textbf{L}ógica \textbf{A}ritimétrica) é um dos componentes mais básicos de um processador. Como o nome já indica, a ULA é responsável por todas as operações lógicas (como OR, AND e Shift) e aritiméticas (como soma e subtração) realizadas em nosso circuito.

De maneira simplificada, a ULA receberá um comando, composto por seletores de operação e bits, e operandos. Na saída temos o resultado da operação desejada.

Como uma ULA opera de maneira concorrente, todas as operações implementadas são ``executadas ao mesmo tempo''. Um Multiplexador é usado para selecionar a operação correta.

\subsection{ULA Proposta}

Neste trabalho prático, a ULA proposta deve possuir as seguintes funções:

\begin{center}
    \begin{tabular}{|c|c|c|c|c|c|}
        \hline
        Função & op\_code & Descrição & z\_out & c\_out & dc\_out\\
        \hline
        XOR & 0000 & XOR Lógico & 1, se res. = 0 & - & -\\
        \hline
        OR & 0001 & OR Lógico & 1, se res. = 0 & - & -\\
        \hline
        AND & 0010 & AND Lógico & 1, se res. = 0 & - & -\\
        \hline
        CLR & 0011 & Limpa & 1 & - & -\\ 
        \hline
        ADD & 0100 & Soma & 1, se res. = 0 & 1 se \textit{carry} & 0 se \textit{carry} no nibble\\
        \hline
        SUB & 0101 & Subtração & 1, se res. = 0 & 0 se \textit{borrow} & 0 se \textit{borrow} no nibble\\
        \hline
        INC & 0110 & Incremento & 1, se res. = 0 & - & - \\
        \hline
        DEC & 0111 & Decremento & 1, se res. = 0 & - & - \\
        \hline
        PASS\_A & 1000 & Passa 'A' & 1, se res. = 0 & - & - \\
        \hline
        PASS\_B & 1001 & Passa 'B' & 1, se res. = 0 & - & - \\
        \hline
        COM & 1010 & Complemento & 1, se res. = 0 & - & - \\
        \hline
        SWAP & 1011 & Permuta \textit{nibbles} & 1, se res. = 0 & - & - \\
        \hline
        BS & 1100 & bit\_sel = 1 & a\_in[bit\_sel] & - & - \\
        \hline
        BC & 1101 & bit\_sel = 0 & a\_in[bit\_sel] & - & - \\
        \hline
        RR & 1110 & Rotação p/ dir. & - &  a\_in[0] & - \\
        \hline
        RL & 1111 & Rotação p/ esq. & - &  a\_in[7] & - \\
        \hline
    \end{tabular}
\end{center}

O sinais de entrada e saída são os seguintes:

\begin{center}
\begin{tabular}{|c|c|c|c|}
    \hline
    Nome & Tamanho & Tipo & Descrição\\
    \hline
    a\_in & 8 bits & \textit{Input} & Entrada de dados A\\
    \hline
    b\_in & 8 bits & \textit{Input} & Entrada de dados B\\
    \hline
    c\_in & 1 bit & \textit{Input} & Entrada de \textit{carry}\\
    \hline
    op\_sel & 4 bits & \textit{Input} & Seletor de operação\\
    \hline
    bit\_sel & 3 bits & \textit{Input} & Seletor de bit\\
    \hline
    r\_out & 8 bits & \textit{Output} & Saída do resultado\\
    \hline
    c\_out & 1 bit & \textit{Output} & Saída de \textit{carry/borrow}\\
    \hline
    dc\_out & 1 bit & \textit{Output} & Saída de \textit{digit carry/borrow}\\
    \hline
    z\_out & 1 bit & \textit{Output} & Saída de zero\\
    \hline
\end{tabular}
\end{center}

\section{Conclusion}
\lipsum[4]

\end{document}